\documentclass[11pt]{article}
\usepackage[margin=1in]{geometry}
\usepackage{xcolor}
\usepackage{hyperref}
\usepackage{enumitem}

% Reviewer comment style
\newcommand{\comment}[1]{\medskip\noindent\textbf{#1}\medskip}
\newcommand{\response}[1]{\noindent #1\medskip}
\newcommand{\done}{\textcolor{blue}{[Done]}}
\newcommand{\pending}{\textcolor{red}{[Pending]}}

\begin{document}

\noindent
\textbf{\Large Response to Reviews}\\[4pt]
\textbf{Manuscript:} USYB-2024-141\\
\textbf{Title:} TreeFlow: Probabilistic Modelling and Automatic Differentiation for Phylogenetics\\
\textbf{Authors:} Swanepoel, Fourment, Ji, Nasif, Suchard, Matsen~IV, Drummond\\[12pt]

\noindent
We thank the Editor-in-Chief, the Associate Editor, and all three reviewers for their thorough and constructive feedback. We have addressed all points raised. Below we describe the changes made; all are indicated in the revised manuscript.

\bigskip
\hrule
\bigskip

%% ============================================================
\section*{Editor-in-Chief (Dr Isabel Sanmart\'in)}
%% ============================================================

\comment{Improvements are needed in the description of the software, particularly the instructions for installation and running the analyses in the tutorial. Both reviewers \#1 and \#2 reported difficulties in running the software.}

\response{\done\ We have resolved all installation and runtime issues reported by Reviewers~1 and~2. TreeFlow now installs cleanly via \texttt{pip install .} in a virtual environment, the tutorial runs to completion, and the Docker image builds successfully. We have also written comprehensive documentation for all associated repositories. A detailed list of software fixes is provided in CHANGES.md in the repository.}

\comment{Reviewer \#3 has provided a detailed report with many useful suggestions\ldots on how to improve the figures and the description of the user interface, as well as the inclusion of additional references.}

\response{\done\ We have addressed all of Reviewer~3's suggestions, including reorganising the manuscript structure, improving figures, adding the requested references, and moving detailed code listings to a supplementary appendix. See our point-by-point responses below.}

\bigskip
\hrule
\bigskip

%% ============================================================
\section*{Reviewer 1}
%% ============================================================

\comment{Command-line documentation does not explain the output produced by executables, does not provide an example of how to run the executable.}

\response{\done\ We have expanded the command-line documentation to include example invocations and descriptions of all output files.}

\comment{Installation error: \texttt{pip install -{}-user .} fails inside a virtual environment.}

\response{\done\ The erroneous \texttt{-{}-user} flag has been removed from the installation instructions.}

\comment{\texttt{treeflow\_vi} fails with missing \texttt{tf\_keras} module.}

\response{\done\ This missing dependency has been added and the CLI now works correctly.}

\comment{Tutorial breaks at \texttt{model\_dist.sample()} with an \texttt{AttributeError}.}

\response{\done\ This was caused by an internal TensorFlow API change. We have replaced the affected code with a stable public API. The tutorial now runs to completion.}

\comment{Docker build fails.}

\response{\done\ The Docker build has been fixed (updated dependencies and Dockerfile syntax).}

\comment{The usability and level of documentation of TreeFlow are still too low to allow the wider community to build upon it.}

\response{\done\ In addition to fixing all reported installation and runtime issues, we have written comprehensive documentation for the TreeFlow, treeflow-paper, treeflow-benchmarks, and phylojax repositories, updated the installation guide, and fixed the API documentation build.}

\bigskip
\hrule
\bigskip

%% ============================================================
\section*{Reviewer 2}
%% ============================================================

\comment{Pip installation: importing treeflow gives a missing \texttt{tf\_keras} error.}

\response{\done\ See response to Reviewer~1. This missing dependency has been added.}

\comment{\texttt{model\_dist.sample()} generates an \texttt{AttributeError}.}

\response{\done\ See response to Reviewer~1. The underlying API incompatibility has been fixed.}

\comment{Docker build fails.}

\response{\done\ See response to Reviewer~1.}

\comment{Overall positive; requests issues be resolved before publication.}

\response{Thank you for the supportive assessment. All reported issues have been resolved and verified.}

\bigskip
\hrule
\bigskip

%% ============================================================
\section*{Reviewer 3 (Fredrik Ronquist)}
%% ============================================================

We are grateful for the detailed and constructive review. We have implemented all suggested changes.

\subsection*{General comments}

\comment{The user interface should be emphasised better in Fig.~1 and introduced first in the text, before the Python API and the architecture of the core.}

\response{\done\ We have reorganised the Software Description section so that the User Interface (CLI and model definition) is presented first, followed by the Developer API. Fig.~1 now uses distinct colours and bolder borders to distinguish the two components.}

\comment{Some text should describe the philosophy of the model specification language, with Figs.~3 and~5 moved to a supplementary appendix with simpler introductory examples.}

\response{\done\ Figs.~3 and~5 have been moved to a new Supplementary Appendix with introductory text on the model specification philosophy. The main text now references Supplementary Figures~S1 and~S2.}

\comment{The Examples section should focus on TreeFlow-related implications rather than biological implications. The Discussion should emphasise TreeFlow advantages over limitations.}

\response{\done\ We have added TreeFlow-focused conclusions to both examples (ease of model extension for carnivores; scalability and YAML-only specification for influenza). The Discussion now leads with TreeFlow's advantages, with limitations reframed alongside forward-looking directions.}

\subsection*{Detailed comments}

\comment{L13--14: Replace RAxML with PhyloBayes; stress Bayesian MCMC focus from the start.}

\response{\done\ Replaced with PhyloBayes (Lartillot et al.\ 2009) and reframed the opening sentence.}

\comment{L19--20: Optionally cite Lakner et al.\ (2008, Syst Bio).}

\response{\done\ Added.}

\comment{L28: Cite H\"ohna et al.\ (2014, Syst Bio) on phylogenetic graphical models.}

\response{\done\ Added.}

\comment{L29--49: ``I like this description of the distinction between probabilistic modelling and probabilistic programming.''}

\response{Thank you.}

\comment{L65: ``of discrete part'' $\to$ ``of the discrete part''.}

\response{\done}

\comment{L78: Cite Senderov et al.\ (2023/2024, TreePPL).}

\response{\done\ Added alongside Ronquist et al.\ (2021).}

\comment{L82: Add ``automatically''; clarify that LinguaPhylo is a tool, not just a language.}

\response{\done}

\comment{L92--94: ``phylogenetic inference'' $\to$ ``phylogenetic problems''.}

\response{\done}

\comment{L98: ``Zhang and IV'' rendering issue.}

\response{\done\ Fixed in all BibTeX entries.}

\comment{L100: ``represention'' $\to$ ``representation''.}

\response{\done}

\comment{L109: ``fixed-topology phylogenetic inference'' $\to$ ``parameter inference in fixed-topology models of a standard form''.}

\response{\done}

\comment{L153: ``probabilsitic'' $\to$ ``probabilistic''.}

\response{\done}

\comment{L161: ``on a marginalizing'' $\to$ ``on marginalizing''.}

\response{\done}

\comment{L171--174: Distinguish Developer API and User Interface more clearly in Fig.~1.}

\response{\done\ See response to the general comment on Fig.~1 above.}

\comment{L189: ``might be need'' $\to$ ``might need''.}

\response{\done}

\comment{L236--239: Difficult to understand; explain ``fixed-topology distribution'' and ``bijectors can be chained'' in simpler terms.}

\response{\done\ We have added a concrete example showing how real-valued parameters are first transformed to ratio representations, then to node heights, and finally combined into a full tree object.}

\comment{L274--276: Sentence needs rewriting.}

\response{\done\ Rewritten for clarity.}

\comment{L324--331: Start the Software Description with this paragraph rather than placing it at the end.}

\response{\done\ Addressed by the reorganisation that places the User Interface description first.}

\comment{Fig.~2: Provide information on potential Monte Carlo error in BEAST estimates.}

\response{\done\ We now show bootstrap density bands around the BEAST MCMC estimates to indicate Monte Carlo sampling error, allowing the reader to assess whether discrepancies between MCMC and VI are within the expected range of Monte Carlo variability.}

\comment{Fig.~3: Move to supplementary material with simpler introductory examples.}

\response{\done\ See general comment above.}

\comment{L369: ``This corrects for some of the posterior approximation error\ldots'' needs a reference and intuitive explanation.}

\response{\done\ Added a citation to Burda et al.\ (2015) and an explanation: the importance weights account for the discrepancy between the variational approximation and the true posterior, reweighting samples to better represent the target distribution.}

\comment{L369--371: Redundant sentence; move numbers to L375 with interpretation.}

\response{\done\ Consolidated the numbers with interpretive context.}

\comment{L384--387: Rewrite conclusion to focus on TreeFlow implications.}

\response{\done\ See general comment on Examples conclusions above.}

\comment{Fig.~5: Move to supplementary with introduction to the model specification language.}

\response{\done\ See general comment above.}

\comment{Fig.~6a: Indicate Monte Carlo error in BEAST estimates.}

\response{\done\ We note that Fig.~6a shows per-lineage kappa estimates from TreeFlow's variational inference, not from BEAST MCMC; there are no BEAST estimates in this figure. Because VI produces independent samples, Monte Carlo error is negligible (approximately $\mathrm{SD}/\sqrt{1000}$ per parameter). We have regenerated this figure with a log-scaled y-axis, which better displays the declining kappa trend across lineages.}

\comment{L428: How much time was spent tuning BEAST? Autotuning? VI convergence diagnostic?}

\response{\done\ We have clarified that BEAST used its built-in autotuning procedure supplemented by manual operator weight adjustment, and that ELBO monitoring was used as the VI convergence diagnostic.}

\comment{L474ff: Write out numbers under ten.}

\response{\done}

\comment{Fig.~7: Reverse legend symbol order.}

\response{\done}

\comment{L529--538: More detail needed on ``structured posterior approximations''.}

\response{\done\ Expanded with concrete examples: block-diagonal covariance structures that capture correlation between clock rate and tree height, and normalizing flows for more flexible approximations.}

\comment{L539--545: Outline how topology inference could be addressed.}

\response{\done\ We describe how subsplit Bayesian networks could be combined with TreeFlow's continuous-parameter VI components to enable joint topology and parameter inference.}

\comment{L562: ``it's'' $\to$ ``its''.}

\response{\done}

\comment{L553ff: Start the Discussion with this text (TreeFlow advantages).}

\response{\done\ See general comment on Discussion reordering above.}

\bigskip
\hrule
\bigskip

\section*{Additional changes}

\begin{itemize}[nosep]
  \item Updated references that have since been published (LinguaPhylo, Ji et al., Fisher et al., Senderov et al.).
  \item Specified BEAST version as 2.7.7 at first analysis mention.
  \item Re-ran BEAST analyses for both datasets to obtain clean wall-clock timing.
\end{itemize}

\end{document}
